\newpage
\chapter{KESIMPULAN DAN SARAN} \label{Bab V}

\section{Kesimpulan} \label{V.Kesimpulan}

Berdasarkan hasil perancangan, implementasi, dan pengujian Sistem Informasi Penerimaan Mahasiswa Baru (PMB) Universitas HKBP Nommensen berbasis web, dapat ditarik beberapa kesimpulan sebagai berikut.\par

\begin{enumerate}
	\item Sistem PMB berbasis web berhasil dibangun menggunakan kerangka kerja Spring Boot 3.3.0 (\textit{back-end}), HTML/Bootstrap 5 (\textit{front-end}), dan MySQL 8.0 sebagai sistem manajemen basis data. Sistem mencakup 22 modul inti yang menangani seluruh alur penerimaan mahasiswa baru, mulai dari pembuatan akun camaba, pembayaran formulir melalui Virtual Account BRIVA, pengisian dan validasi formulir pendaftaran, pelaksanaan ujian seleksi, pengumuman kelulusan, proses daftar ulang, hingga ekspor dan pelaporan data. Arsitektur \textit{stateless} berbasis JWT memastikan keamanan autentikasi multi-peran (Admin Pusat, Admin Validasi, dan Camaba) tanpa bergantung pada manajemen sesi \textit{server-side}.
	
	\item Metode \textit{Evolutionary Prototyping} terbukti efektif sebagai pendekatan pengembangan sistem PMB yang bersifat dinamis. Pengembangan berlangsung dalam empat iterasi: Iterasi 1 menyelesaikan 10 item perbaikan dari evaluasi prototipe awal (termasuk sistem cicilan enam tahap, fitur ekspor multi-format, dan manajemen \textit{system links}); Iterasi 2 memperbaiki mekanisme cicilan berurutan, threshold \textit{autocomplete} sekolah, dan tampilan ekspor formulir; Iterasi 3 menambahkan manajemen data SMA oleh admin, unduhan bundel dokumen daftar ulang dalam format ZIP, dan penonaktifan akun pengguna; serta Iterasi 4 menyelesaikan finalisasi \textit{real-time preview} jadwal cicilan, penanganan kesalahan ekspor ZIP, dan optimisasi kueri pencarian sekolah menggunakan indeks \texttt{FULLTEXT}. Seluruh kebutuhan fungsional PMB01-F002 hingga PMB01-F047 terpenuhi pada akhir Iterasi 4.
	
	\item Hasil \textit{Black-Box Testing} yang dilakukan dengan 21 kasus uji (TC001--TC021) mencakup 40 skenario pengujian---termasuk skenario positif dan negatif---menunjukkan bahwa seluruh kasus uji berstatus \textit{Pass}. Sistem berhasil menangani validasi masukan, pengendalian akses berbasis peran, dan pengiriman notifikasi email secara otomatis. Hasil \textit{White-Box Testing} pada 15 jalur eksekusi di empat komponen kritis (\texttt{JwtTokenProvider}, \texttt{ValidasiService}, \texttt{CicilanService}, dan \texttt{AdminController}) mencapai \textit{branch coverage} 100\%, memastikan setiap kondisi percabangan dan penanganan pengecualian berjalan sesuai spesifikasi.
	
	\item Hasil pengujian \textit{usability} menggunakan \textit{System Usability Scale} (SUS) yang melibatkan 20 responden (15 camaba dan 5 admin) menghasilkan skor rata-rata \textbf{84{,}25}. Berdasarkan skala adjektif SUS menurut Bangor et al.\ (2009), skor ini masuk dalam kategori \textit{Excellent} dengan \textit{Grade A-} dan status \textit{Acceptable}, yang mengindikasikan bahwa sistem memiliki tingkat \textit{usability} yang tinggi dan dapat diterima dengan sangat baik oleh pengguna akhir. Sebanyak 90\% responden (18 dari 20) memberikan penilaian dalam kategori \textit{Excellent}, menunjukkan konsistensi kualitas antarmuka di berbagai kelompok pengguna.
\end{enumerate}

\section{Saran} \label{V.Saran}

Berdasarkan hasil penelitian dan keterbatasan yang ditemukan selama pengembangan, terdapat beberapa saran untuk pengembangan sistem PMB ke depannya.\par

\begin{enumerate}
	\item \textbf{Integrasi dengan Sistem Akademik (SIAKAD) UHN.} Sistem PMB saat ini beroperasi sebagai sistem terpisah dari sistem akademik yang sudah ada di Universitas HKBP Nommensen. Pengembangan selanjutnya disarankan untuk membangun antarmuka integrasi (\textit{API bridge}) antara sistem PMB dan SIAKAD UHN, sehingga data mahasiswa baru yang berhasil daftar ulang dapat secara otomatis masuk ke sistem administrasi akademik tanpa entri data ulang. Integrasi ini akan menghilangkan potensi ketidakkonsistenan data antara dua sistem.
	
	\item \textbf{Pengembangan Aplikasi \textit{Mobile}.} Hasil pengujian SUS menunjukkan bahwa sebagian camaba mengalami hambatan dalam mengakses beberapa halaman formulir yang panjang melalui perangkat \textit{mobile} meskipun tampilan sudah responsif. Pengembangan aplikasi \textit{mobile} native (Android/iOS) atau \textit{Progressive Web App} (PWA) dapat memberikan pengalaman pengguna yang lebih baik, terutama untuk alur penting seperti pengecekan status pendaftaran, notifikasi \textit{push}, dan unggah dokumen langsung dari kamera perangkat.
	
	\item \textbf{Penambahan Fitur Analitik dan Pelaporan Manajemen.} Sistem saat ini menyediakan ekspor data dalam format CSV dan JSON, namun belum memiliki dasbor analitik visual. Pengembangan modul analitik yang menampilkan tren pendaftar per program studi, tingkat konversi per tahap alur (registrasi $\rightarrow$ bayar $\rightarrow$ daftar $\rightarrow$ lulus $\rightarrow$ daftar ulang), dan perbandingan antar gelombang akan memberikan nilai tambah yang signifikan bagi manajemen universitas dalam pengambilan keputusan penerimaan mahasiswa.
	
	\item \textbf{Peningkatan Keamanan Sistem.} Meskipun sistem sudah menerapkan autentikasi JWT dan enkripsi \textit{password} BCrypt, beberapa aspek keamanan dapat ditingkatkan untuk lingkungan produksi: (a) penerapan \textit{rate limiting} pada \textit{endpoint} login dan pengiriman email untuk mencegah serangan \textit{brute-force} dan spam, (b) implementasi autentikasi dua faktor (2FA) untuk akun Admin Pusat mengingat hak aksesnya yang luas, serta (c) audit log yang mencatat seluruh aksi kritis (perubahan data camaba, penetapan hasil akhir, penonaktifan akun) beserta identitas pengguna dan \textit{timestamp}-nya untuk keperluan audit dan akuntabilitas.
\end{enumerate}